\chapter{为什么「预测下一个词」能产生智能？}
\label{ch:intro}
%% ============================================================

\section{一个看似简单的任务}

假设你看到这样一个句子的开头：「法国的首都是\_\_\_\_」。你会毫不犹豫地填入「巴黎」。
这个填空题看起来极其简单，但请思考一下：要正确完成这个填空，你需要具备什么？
你需要知道法国是一个国家，首都是一个政治概念，以及巴黎与法国之间的具体关系。
换句话说，你需要拥有\textbf{世界知识}。

再看一个稍微复杂的例子：「她在得知祖母去世后感到\_\_\_\_」。
要预测这里的词，你需要理解人类情感——失去亲人通常会带来悲伤。
这不仅仅是语言知识，更是对人类心理和社会关系的理解。

再来一个：「如果所有的猫都是动物，而 Tom 是一只猫，那么 Tom 是\_\_\_\_」。
这里需要的是逻辑推理能力，三段论的基本形式。

还有一个更微妙的例子：「他把钥匙放在桌子上，然后走进了厨房。
几分钟后他回来拿\_\_\_\_」。要预测「钥匙」，
模型需要维持跨越多个句子的\textbf{工作记忆}：
记住「钥匙」这个对象及其位置，
并在上下文发生变化后依然能够追溯引用它。

最后一个：「$\int_0^1 x^2 \, dx = $\_\_\_\_」。
这里需要的是纯粹的数学计算能力：
理解积分符号的含义、掌握幂函数积分的规则、
执行代入求值。

这就是大语言模型（Large Language Model, LLM）训练的核心目标：
\textbf{准确预测下一个词}。
这个目标看似简单，实则蕴含着一个深刻的洞察：
要在所有可能的上下文中都准确预测下一个词，
模型必须隐式地学会语法、语义、逻辑推理、世界知识、情感理解，
乃至数学和编程。

\begin{corebox}[预测下一个词的数学形式]
给定前 $t$ 个 token $x_1, x_2, \ldots, x_t$，
模型学习条件概率分布：
\[
  P(x_{t+1} \mid x_1, x_2, \ldots, x_t)
\]
训练目标是最大化整个训练语料库上所有位置的对数似然：
\[
  \mathcal{L} = -\sum_{t=1}^{T} \log P(x_t \mid x_1, \ldots, x_{t-1})
\]
这个看似简单的损失函数，驱动了从 GPT-1 到 GPT-5 的全部进化。
\end{corebox}

\section{预测下一个词：一种全新的学习范式}

为什么 next-token prediction 如此特别？要理解这一点，我们需要将它与此前的机器学习范式做深入对比。

\subsection{从特征工程到监督分类}

传统机器学习的核心循环是：人类专家分析问题 $\to$ 设计特征 $\to$ 训练分类器 $\to$ 在测试集上评估。这个范式在本质上受限于人类专家的认知，你能提取的特征决定了模型能力的上限。一位垃圾邮件检测工程师可能设计出「是否包含'免费'关键词」「发件人域名信誉度」等特征，但他永远无法穷举所有垃圾邮件的模式。

监督分类（supervised classification）将这个范式推进了一步：给定输入 $x$ 和标签 $y$，学习 $P(y \mid x)$。深度学习免去了手工特征设计的苦差事，神经网络自动从原始输入中学习特征。但监督学习有一个根本性的瓶颈：\textbf{标注数据}。ImageNet 花费了数年时间，动用了亚马逊 Mechanical Turk 上的 49,000 名标注员，才标注了 1400 万张图片~\cite{deng2009imagenet}。自然语言的标注成本更高——判断一段文本的情感倾向、提取其中的实体关系、评估翻译质量，每一项都需要训练有素的标注员。

\subsection{BERT 的掩码预测：一个重要但受限的中间步骤}

BERT~\cite{devlin2019bert} 提出了一种聪明的自监督方法：随机遮盖输入中 15\% 的 token，让模型预测被遮盖的部分。这就是 Masked Language Model（MLM）。MLM 的训练数据是免费的，任何文本都可以随机遮盖，无需人工标注。

但 MLM 有几个深层局限。首先，它是\textbf{双向的}，模型同时看到被遮盖位置的左右上下文。这意味着 BERT 不是一个生成模型：它不能自然地逐词生成文本，因为它在训练时就没有学过「只看到前文、预测后续」这种自回归模式。其次，MLM 一次只预测 15\% 的 token，每个训练样本的信号密度远低于自回归模型（后者对每个位置都产生梯度信号）。第三，也是最微妙的一点：遮盖是\textbf{独立随机}的，BERT 假设被遮盖的 token 之间互相独立，但自然语言中相邻的词往往强相关~\cite{yang2019xlnet}。

\subsection{Next-Token Prediction 的独特力量}

相比之下，next-token prediction 具备几个关键优势，使其成为通向通用智能的最有效路径。

\textbf{训练信号的密度}。自回归模型对序列中的\textbf{每一个}位置都计算损失，没有任何 token 被「浪费」。一个长度为 $T$ 的序列提供 $T-1$ 个训练信号，而 MLM 只提供约 $0.15T$ 个。这意味着自回归训练的数据效率比 MLM 高出约 6--7 倍。当训练数据的规模达到万亿 token 时，这个差距是决定性的。

\textbf{生成能力的天然获得}。自回归模型在训练时就在学习「给定前文，预测下一个词」，这与生成任务完全一致。训练完成后，只需从学到的概率分布中采样，模型就能生成连贯的文本。BERT 要做到这一点需要复杂的额外工程。

\textbf{隐式的多任务学习}。这是最深刻的一点。预测下一个词\textbf{天然地}蕴含了几乎所有的 NLP 任务。要预测一段翻译的下一个词，模型需要掌握翻译能力。要预测一段推理过程的下一步，模型需要掌握逻辑推理。要预测一段代码的下一行，模型需要理解编程。Radford 等人在 GPT-2 论文中精辟地指出：语言建模可以被视为一种无监督的多任务学习~\cite{radford2019gpt2}。当训练语料包含了人类知识的方方面面时，next-token prediction 就隐式地学习了所有这些能力。

\textbf{零样本泛化}。这是上述多任务学习的直接推论。如果模型在训练中见过「将英语翻译成法语」的例子，它就能在推理时执行翻译，无需任何额外训练。GPT-3 展示的 in-context learning~\cite{brown2020gpt3} 正是这种能力的具体体现：通过在提示中给出几个示例，模型就能「学会」新任务。这从根本上改变了 AI 系统的使用方式——从「每个任务训练一个模型」变成「一个模型解决所有任务」。

\textbf{训练数据的免费性}。这是最容易被忽略但最重要的优势。监督学习需要昂贵的标注数据；BERT 的 MLM 虽然自监督，但每个训练样本只利用了 15\% 的 token。而 next-token prediction 的训练数据是\textbf{完全免费的}，任何文本都是有效的训练样本，每个 token 都提供梯度信号。互联网上存在数十万亿 token 的文本数据，且每天都在增长。这意味着 next-token prediction 的训练数据在理论上是\textbf{无限的}——唯一的瓶颈是计算量和数据质量，而非数据数量本身。（当然，如第~\ref{sec:scaling-laws}节所述，高质量数据正在接近耗尽，但这是质量的瓶颈，不是数量的瓶颈。）

这个优势的经济含义是巨大的。
ImageNet 的标注成本约为 50 万美元（2009--2012 年），
涵盖 1400 万张图片。
按照同样的标注质量标准，
标注一个 GPT-3 级别的文本数据集
（300B tokens，约 200 万篇长文档）
可能需要数亿美元和数年时间。
而 next-token prediction 将这个成本降为\textbf{零}：
数据已经标注好了（下一个词就是标签），
你只需要收集和清洗数据本身。
这一经济优势是 LLM 范式能够扩展到万亿 token 规模的根本原因。

\begin{keybox}[两种范式的本质区别]
\begin{itemize}[noitemsep]
  \item \textbf{监督分类}：人类定义任务 $\to$ 收集标注数据 $\to$ 训练专用模型。能力上限由标注数据的质量和数量决定。
  \item \textbf{BERT-MLM}：自监督预训练 $\to$ 每个下游任务微调一个模型。消除了标注瓶颈，但仍需要任务特定的适配。
  \item \textbf{Next-token prediction}：在海量文本上预训练 $\to$ 单个模型通过提示完成任意任务。消除了任务定义的瓶颈。
\end{itemize}
范式的演进方向是清晰的：\textbf{人类干预越来越少，模型通用性越来越强}。
\end{keybox}

\subsection{为什么自回归胜出了？}

一个常被问及但不易回答的问题：为什么是 decoder-only 的自回归架构（GPT 路线）胜出，而不是 encoder-only（BERT 路线）或 encoder-decoder（T5 路线）？

答案不在数学优雅性上——三种架构在理论表达能力上是等价的。差异来自\textbf{实践中的规模效应}。

Decoder-only 架构有一个工程上的巨大优势：\textbf{所有 token 共享同一个前向计算路径}。不需要像 encoder-decoder 那样维护两个不同的模块，也不需要处理 encoder 和 decoder 之间的 cross-attention。这使得训练和推理的工程实现更简单，分布式并行化更高效。当模型规模从 1B 增长到 1000B 时，工程简洁性的优势被放大，更少的设计选择意味着更少的 bug 和更容易的调优。

另一个常被忽视的优势是\textbf{KV 缓存}的高效性。
在自回归生成中，decoder-only 架构可以缓存之前所有 token 的 Key 和 Value，
生成新 token 时只需要计算当前 token 的 Query 与已缓存 KV 的注意力：
这使得每生成一个 token 的边际成本远低于重新计算整个序列。
Encoder-decoder 架构在 decoder 部分也有 KV 缓存，
但 cross-attention 部分还需要额外缓存 encoder 的输出：
内存开销更大。
在高并发的推理服务中
（如 ChatGPT 同时服务数百万用户），
KV 缓存的效率直接决定了服务的吞吐量和成本。

还有一个理论上的优势。
Decoder-only 架构的因果掩码（causal mask）
确保了模型在训练时只能看到当前位置之前的 token：
这与推理时的生成过程完全一致。
BERT 的双向注意力在训练时可以看到完整的上下文，
但在生成时无法做到这一点：
存在\textbf{训练-推理不一致}（train-inference mismatch）的问题。
XLNet~\cite{yang2019xlnet} 试图用排列语言模型来解决这个问题，
但增加的复杂性使其未能广泛采用。

更深层的原因是：自回归训练天然支持无限长的上下文。模型可以处理任意长度的输入，只要硬件内存允许。而 BERT 的双向注意力在处理长序列时面临更大的效率挑战。2024--2025 年，随着上下文窗口从 8K 扩展到 1M 乃至 10M token，自回归架构的这一优势变得愈发关键。

\subsection{Next-Token Prediction 的理论深度}

Next-token prediction 的力量还有更深的理论根源。2024 年，Hahn 和 Goyal~\cite{hahn2024theory} 的理论分析表明，一个足够大的 Transformer 在训练 next-token prediction 时，它不仅仅在学习表面的 n-gram 统计，它实际上在学习一种\textbf{隐式的世界模型}。

直觉如下。考虑训练语料中关于物理世界的描述：「球滚到了桌子边缘，然后\_\_\_\_」。要正确预测「掉了下去」，模型需要建立一个关于重力、桌子边缘和球的运动的内部表示。如果训练语料足够丰富，模型面对的预测任务涵盖了物理世界的方方面面，模型被迫在其参数中编码一个越来越精确的世界模型。

Li 等人~\cite{li2023othello} 在 Othello 棋盘游戏上提供了直接证据。他们训练了一个 GPT 模型来预测 Othello 棋步序列中的下一步。模型从未被告知棋盘的存在，它只看到了一系列棋步记录。但研究者发现，模型的内部表示中\textbf{涌现出了棋盘状态的线性表示}，通过简单的线性探针（linear probe），可以从模型的中间激活值中读出当前棋盘上每个位置的状态（黑子、白子或空）。

这个发现的含义深远：next-token prediction 不仅是在学习序列的统计模式，而是在学习\textbf{生成这些序列的底层过程}。对于自然语言来说，这个底层过程就是人类的思维、世界知识和推理链。当然，模型学到的「世界模型」可能是不完整的、有偏差的、甚至是错误的，但它确实存在，而且随着规模的增加在变得更加精确。

\subsection{局限性：Next-Token Prediction 不是万能的}

在强调 next-token prediction 的力量之后，有必要诚实地面对它的局限。

\textbf{概率口吃}（probability smearing）。语言模型学到的是所有合理延续的加权平均。当一个上下文有多个同样合理的延续时，模型的预测会「模糊」，它不会坚定地选择一个方向，而是在多个选项之间摇摆。这在需要一致性的长文本生成中尤为明显（例如，模型可能在一个故事中给主角取了不同的名字）。这个问题的根源在于：next-token prediction 是在\textbf{所有合理文本的分布}上训练的，而不是在单个连贯叙事上训练的。模型学到的不是「如何讲一个特定的故事」，而是「所有可能的故事的混合」。后训练阶段的 RLHF/DPO 可以部分缓解这个问题，但无法从根本上消除它。

\textbf{因果推理的缺失}。语言模型学到的是相关性，不是因果性。「冰淇淋销量上升」和「溺水事故增加」在训练数据中是相关的（因为都与夏天相关），但模型可能无法正确区分相关和因果。Pearl~\cite{pearl2009causality} 的因果推理理论指出，纯粹的观测数据（无干预）从原理上就不足以学习因果关系，而 next-token prediction 的训练数据正是纯观测的。

\textbf{推理深度的限制}。Transformer 的每一层执行一次「计算步骤」。一个 $L$ 层的 Transformer 在处理每个 token 时，最多执行 $L$ 步计算。某些推理任务（如需要 $O(n)$ 步的算法问题）可能超出了固定深度 Transformer 的能力~\cite{merrill2023expressive}。这正是 Chain-of-Thought 和 test-time compute 的价值所在，通过生成中间步骤，模型将受限的单步计算深度扩展为不受限的多步推理链。

\section{压缩即理解}

信息论的创始人 Claude Shannon 在 1948 年就揭示了这个联系~\cite{shannon1948}：
\textbf{预测和压缩是同一件事的两面}。
如果你能完美预测一段文字的下一个字符，
你就可以用最少的比特来编码它，这就是最优压缩。
反过来，最优压缩器必然是最优预测器。

\subsection{信息熵与交叉熵损失}

Shannon 定义了信息熵（entropy）来度量信息源的不确定性：
\begin{equation}
  H(X) = -\sum_{x} p(x) \log_2 p(x)
\end{equation}
其中 $p(x)$ 是符号 $x$ 的真实概率。
熵越高，信息源越不可预测，需要的编码位数越多。

Shannon 在 1950 年的论文《Prediction and Entropy of Printed English》~\cite{shannon1950}中通过实验估算，英语的熵大约在 \textbf{0.6--1.3 bits/字符}之间。
这意味着如果你有一个完美的英语预测器，
每个字符平均只需要约 1 bit 来编码：
远少于使用 ASCII 编码的 7 bits/字符。
换句话说，英语文本有大约 75--85\% 的\textbf{冗余度}，
而这些冗余恰恰来自语言的规律性——语法、语义、常识，
正是模型需要学习的知识。

当我们用模型 $q$ 来近似真实分布 $p$ 时，
编码效率由\textbf{交叉熵}（cross-entropy）衡量：
\begin{equation}
  H(p, q) = -\sum_{x} p(x) \log_2 q(x) \geq H(p)
\end{equation}
交叉熵永远大于等于真实熵——差值就是 KL 散度，
衡量模型 $q$ 与真实分布 $p$ 之间的差距。
LLM 训练时最小化的损失函数正是交叉熵：
\textbf{训练 LLM 本质上是在训练一个压缩器}。

\subsection{困惑度：压缩效率的直觉度量}

困惑度（perplexity）是交叉熵的指数形式：
\begin{equation}
  \mathrm{PPL} = 2^{H(p, q)}
\end{equation}
直觉上，困惑度表示模型在预测下一个 token 时的「平均选项数」。
PPL $= 10$ 意味着模型在每一步平均面临 10 个等可能的选择。
如果模型对语言有更好的理解（更好的压缩），选项数更少，PPL 更低。

一个具体的数字感受：GPT-2（2019, 1.5B）在 WikiText-103 上的 zero-shot PPL 约为 37.5，
而 GPT-3（2020）约为 20，GPT-4 级别的模型据推测可能低于 10（OpenAI 未公开具体数字）。
PPL 的下降直接对应着模型对语言理解的提升。

\subsection{Sutskever 的压缩哲学}

OpenAI 的联合创始人 Ilya Sutskever 将这个思想推向了极致。
他在多次演讲中强调：「预测下一个词的过程，就是压缩人类知识的过程。
而足够好的压缩，本质上就是理解。」

这个论断并非修辞。考虑一段关于量子力学的论文：
如果模型能准确预测下一个词，
它必须在某种程度上「理解」量子力学的概念和推理链。
一个仅仅记忆了表面统计模式的模型
（比如「量子」后面常跟「力学」）
在面对新的推理组合时会失败。
只有真正建立了概念之间因果关系的模型
才能在新颖的上下文中做出正确预测。

Deepmind 在 2023 年发表的研究~\cite{deletang2024language}
将这一直觉形式化：
他们证明足够强大的语言模型本质上是通用的数据压缩器，
因此 next-token prediction 在理论上蕴含了强大的表示能力。
当然，「足够强大」和「实际训练出来」之间还有巨大的鸿沟：
但这至少为「压缩即理解」提供了理论基础。

这就是整个 LLM 训练范式的哲学基础：
我们不直接教模型「什么是知识」，
而是通过让它预测海量文本中的下一个词，
迫使它自发地组织和内化人类知识。

\subsection{压缩视角的实践启示}

压缩-理解等价性不仅是哲学思辨，它还有直接的实践启示。

\textbf{训练数据的价值取决于其不可预测性}。如果一段文本对模型来说是高度可预测的（即模型已经「压缩」了这段文本中的知识），那么继续在这段文本上训练带来的信息增益很小。反过来，如果一段文本让模型困惑（高 perplexity），说明模型尚未学会这段文本中的模式，这就是有价值的训练信号。这个洞察催生了\textbf{课程学习}（curriculum learning）策略：先用简单文本训练模型（让它快速建立基础知识），再逐步引入更困难的文本（让它学习更精细的模式）。FineWeb~\cite{penedo2024fineweb} 的数据质量评分本质上就是在度量数据的「可压缩性」——高质量数据包含更多不可预测的、有信息价值的模式。

\textbf{模型的能力可以用压缩率来度量}。如果我们用模型 A 压缩一段文本得到 $x$ bits，用模型 B 压缩相同文本得到 $y$ bits，且 $x < y$，那么模型 A 对这段文本有更好的「理解」——无论这段文本是英语散文、Python 代码还是数学证明。这提供了一种统一的、不依赖特定任务的模型评估方式。实际上，perplexity 就是这种评估方式的直接实现，它度量的正是模型作为压缩器的效率。

\textbf{Kolmogorov 复杂度的联系}。从更深层的理论视角看，完美的语言模型等价于一个计算 Kolmogorov 复杂度（即字符串的最短描述长度）的算法。当然，Kolmogorov 复杂度是不可计算的，我们只能逼近它，不能精确计算它。但 LLM 可以被视为一种实用的近似：它通过在海量数据上训练，学会了用参数来编码数据中的规律，从而提供了一种\textbf{实用的通用压缩器}。这个理论联系解释了为什么更大的模型更「聪明」，它们有更多的参数来编码更精细的规律，从而实现更好的压缩。

\begin{keybox}[压缩--预测--理解 三位一体]
\begin{itemize}[noitemsep]
  \item \textbf{压缩}：用更少的 bit 表示信息 $\Leftrightarrow$ 找到数据中的规律
  \item \textbf{预测}：准确预测下一个 token $\Leftrightarrow$ 理解上下文的结构
  \item \textbf{理解}：建立概念之间的关系 $\Leftrightarrow$ 实现高效压缩
  \item 三者在数学上等价：最小化交叉熵 $=$ 最大化压缩率 $=$ 最优预测
\end{itemize}
\end{keybox}

\section{苦涩的教训}

强化学习之父 Rich Sutton 在 2019 年写了一篇影响深远的短文，
题为「苦涩的教训」（The Bitter Lesson）~\cite{sutton2019bitterlesson}。
他的核心观点是：在 AI 研究的 70 年历史中，
\textbf{利用计算能力的通用方法，最终总是击败人工设计的巧妙方法}。

\subsection{历史佐证}

\textbf{国际象棋}。1997 年，IBM 的 Deep Blue 击败了世界冠军 Kasparov。
Deep Blue 的核心并非对国际象棋有多深刻的「理解」，
而是每秒评估 2 亿个棋局的暴力搜索能力，
辅以相对简单的评估函数。
棋手们花了几十年精心设计的开局库和战术知识库，
最终败给了摩尔定律驱动的算力增长。

\textbf{围棋}。围棋的搜索空间比国际象棋大数十个数量级
（$\sim 10^{170}$ 种合法局面），暴力搜索不再可行。
2016 年，DeepMind 的 AlphaGo 通过深度神经网络 + 蒙特卡洛树搜索击败了李世乭。
更惊人的是 AlphaGo Zero（2017）：
它完全不使用人类棋谱，仅通过自我对弈学习，
40 小时内就超越了所有使用人类知识训练的版本。
这是对 Bitter Lesson 最戏剧化的诠释：
\textbf{人类知识不仅不必要，甚至是一种限制}。

\textbf{语音识别}。从 1970 年代到 2010 年代，
语音识别的主流方法是隐马尔可夫模型（HMM）加上
精心设计的声学特征（MFCC）和发音词典。
这套系统凝聚了数十年的语言学知识。
2012 年后，端到端的深度学习方法
（先是 DNN-HMM 混合系统，后是纯端到端的 CTC 和 attention 模型）
逐步取代了传统方法。
到 2023 年，OpenAI 的 Whisper 模型
在 68 万小时的弱标注数据上训练，
无需任何语言学先验知识就达到了人类水平的识别精度。

\textbf{机器翻译}。1990 年代，基于规则的机器翻译
（如 Systran）需要语言学家手工编写数千条转换规则。
2000 年代，统计机器翻译（如 Moses）用双语平行语料自动学习翻译概率。
2014 年后，神经机器翻译（seq2seq + attention）再次颠覆了整个领域。
到 2017 年 Transformer 出现时，
机器翻译已经完全是一个端到端学习问题，
不需要任何语言学知识。

\textbf{计算机视觉}。2012 年之前，
图像分类的主流方法是手工设计特征（SIFT、HOG、LBP）
加上经典分类器（SVM）。
AlexNet 在 ImageNet 上的胜利标志着这套范式的终结。
此后，更深的网络（VGG、ResNet）和更多的数据
不断推高性能，手工特征被彻底淘汰。
值得注意的是，这个转变的规模并不大：
AlexNet 只有 6000 万参数，训练数据（ImageNet）只有 120 万张图片。
关键不是绝对规模，而是从「人工设计特征」到「学习特征」的范式转换。
LLM 的故事完美地重演了这个模式，只是在更大的规模上。

\subsection{Bitter Lesson 的量化证据}

Bitter Lesson 不仅仅是一种定性的哲学观察：
它有可以量化的证据支持。

考虑自然语言处理中「情感分析」这个经典任务。
2010 年代，最先进的方法是手工设计的特征
（如情感词典匹配、否定检测、讽刺识别规则）
加上经典的机器学习分类器。
一个典型的系统可能有 50--100 条人工规则，
需要一个 NLP 专家花费数月来精心调优，
在标准数据集上达到约 85\% 的准确率。

2018 年，BERT 的 fine-tuning 以零条人工规则达到了 94\% 的准确率。
2020 年，GPT-3 的 few-shot 以零条人工规则达到了 92\% 的准确率：
甚至不需要针对情感分析进行微调。
2023 年，GPT-4 的 zero-shot 达到了约 95\% 的准确率：
只需要一句「判断以下文本的情感倾向」的提示。

这个案例展示了 Bitter Lesson 的三个阶段：
(1) 人工规则在数据稀缺时有价值（2010 年代）；
(2) 当模型规模和数据量足够大时，人工规则被端到端学习取代（2018--2020）；
(3) 最终，通用模型甚至不需要针对特定任务训练（2023+）。
每次转变都是「计算 + 数据」对「人类巧思」的胜利。

\subsection{反面论证：人类知识的价值}

Bitter Lesson 并非没有争议。
一个重要的反面论证是：
\textbf{正确的归纳偏置（inductive bias）能大幅提升学习效率}。

卷积神经网络（CNN）内置了平移不变性和局部性的先验：
这是人类对视觉世界的知识编码。
正是这个先验使得 CNN 在相同数据量下
远优于全连接网络。
Transformer 的注意力机制本身也是一种归纳偏置：
它假设序列中的关系可以通过内容相似度来发现。

更广泛地说，
选择正确的\textbf{架构}（CNN 用于视觉、Transformer 用于序列）
就是一种人类知识的注入。
Bitter Lesson 的真正含义可能不是
「人类知识无用」，而是
「\textbf{手工编码的领域规则}不如\textbf{通用学习算法 + 计算规模}，
但\textbf{结构性先验}（如架构选择）仍然至关重要」。

\subsection{2025 年的新苦涩教训}

LLM 的历史本身就是一连串的苦涩教训。

\textbf{手工规则 vs 端到端学习}。早期的聊天机器人（如微软的 Xiaoice）使用精心设计的对话管理规则、情感分析模块和知识图谱。这些系统凝聚了数百名工程师多年的努力。ChatGPT 用一个端到端训练的语言模型，没有显式的对话管理，没有知识图谱，没有情感分析模块——就在一夜之间使它们全部过时。

\textbf{NLP pipeline vs 通用模型}。在 GPT-3 之前，解决一个 NLP 问题需要一条复杂的 pipeline：分词 $\to$ 词性标注 $\to$ 命名实体识别 $\to$ 句法分析 $\to$ 语义理解 $\to$ 任务特定的推理模块。每个组件都需要独立训练和维护，错误会在 pipeline 中级联放大。一个足够大的语言模型将所有这些步骤融合为单次前向传播，更简单、更可靠、效果更好。

\textbf{搜索增强 vs 参数化知识}。2022 年，RAG（Retrieval-Augmented Generation）被视为解决 LLM 幻觉的关键技术。但到 2025 年，随着模型训练数据量达到数十万亿 token，参数化知识的覆盖面和准确度大幅提升。对于大多数常见问题，GPT-5 和 Claude Opus 4.6 不需要外部检索就能给出准确答案。RAG 仍然在需要实时信息的场景中有价值，但「先搜索再回答」的范式正在被「参数记忆 + 按需搜索」的混合模式取代。

\subsection{LLM 作为 Bitter Lesson 的极致体现}

LLM 的成功正是这个教训的最新、最戏剧化的例证。
我们没有手动编程语法规则、知识库或逻辑推理引擎，
而是选择了最简单的目标函数（预测下一个词），
然后用数万亿的文本和数千块 GPU 进行训练。
结果是：模型涌现出了令人惊叹的能力。

但值得注意的是，Transformer 架构本身
——注意力机制、位置编码、残差连接：
这些设计并非「无知识」的，
而是凝聚了深度学习社区多年的经验和直觉。
从这个角度看，LLM 的成功是
Bitter Lesson 和人类巧思的\textbf{结合}，
而非简单的「规模碾压一切」。

2025--2026 年的效率革命为 Bitter Lesson 增添了一个新维度：
不仅是「计算 + 通用方法胜过巧妙设计」，
还是「\textbf{更聪明的计算利用方式}胜过更大的计算量」。
DeepSeek 用 560 万美元达到 GPT-4 水平：
不是因为它用了更多的 GPU，
而是因为它用了更聪明的架构（MoE + MLA）、
更高效的训练方法（FP8 + 分布式优化）和更好的数据工程。
也许 Bitter Lesson 的真正含义应该更新为：
「\textbf{通用学习算法 + 高效地利用一切可用的计算}，
最终总是胜出。」

\section{涌现：大就是不同}

LLM 的一个最引人入胜的现象是\textbf{涌现}（emergence）。
当模型规模超过某个阈值时，
会突然展现出在较小规模上完全不存在的能力。

\subsection{经典涌现案例}

例如，思维链推理（Chain-of-Thought, CoT）~\cite{wei2022cot}：
小模型在被提示「让我们一步一步思考」时不会产生任何有意义的改进，
但当模型参数超过约 60B 时，
这个简单的提示会带来巨大的推理能力提升。
Wei 等人在 2022 年的研究表明，
CoT 推理能力在模型规模低于 $\sim$10B 时几乎不存在，
在 10--100B 之间开始出现，
在 100B 以上则表现出稳定的强大效果。

算术能力的涌现更为戏剧化。
GPT-3（175B）能做三位数加法但经常出错，
GPT-4 则能可靠地完成多步骤数学推理。
类比推理、代码生成、多语言翻译：
这些能力都不是线性增长的，
而是在模型跨越某个规模门槛时「突然出现」。

另一个引人注目的涌现案例是\textbf{多语言翻译的零样本泛化}。
PaLM（540B）在仅以英语为主的数据上训练后，
展现出了将法语翻译为德语的能力：
即使训练数据中几乎没有法-德平行语料~\cite{chowdhery2022palm}。
模型似乎在大规模训练中建立了一个
「语言无关」的语义空间，
使得它可以在从未明确学习过的语言对之间进行翻译。
这种能力在小模型中完全不存在：
它需要足够大的表示空间来同时编码多种语言的语义结构。

\textbf{指令跟随}本身也是一种涌现能力。
小模型（<10B）在收到「请用一句话概括以下文章」这样的指令时，
通常会忽略指令，直接继续生成文本（因为这是它在预训练中学到的行为）。
但当模型规模增大到一定程度后，
它开始能够「理解」指令并按要求行事：
即使没有经过 SFT（监督微调）。
这个观察表明，指令跟随能力的\textbf{种子}
在预训练阶段就已经被埋下了：
SFT 只是将这个能力「激活」和「强化」，
而非从零创造。

这种涌现现象在物理学中有一个贴切的类比：
\textbf{相变}（phase transition）。
水在 99°C 和 100°C 之间发生的不是渐进变化，而是质变：
从液态到气态。
类似地，语言模型在某个参数规模或训练数据量的临界点上，
可能经历从「不具备某能力」到「具备该能力」的跃迁。
这个类比虽然直觉上吸引人，但严格性有限：
神经网络并非热力学系统，
「相变」是否有精确的数学对应仍是开放问题。

不过，2024--2025 年的研究在形式化这个类比上取得了一些进展。
Michaud 等人发现，在简化的语言模型中，
特定能力的获得确实可以用临界指数来描述：
这表明神经网络的学习动力学可能确实存在
与统计物理中的相变类似的数学结构。
但这些结果目前仅限于高度简化的设定，
能否推广到 GPT 级别的大模型仍不确定。

\subsection{涌现是真实的还是度量假象？}

2023 年，Schaeffer 等人~\cite{schaeffer2023mirage}
发表了一篇引人深思的论文，题为
「Are Emergent Abilities of Large Language Models a Mirage?」。
他们的核心论点是：
许多所谓的涌现可能是\textbf{评估方式的假象}。

具体来说，当使用\textbf{离散}的评估指标
（如准确率，答案要么全对要么全错）时，
能力增长的曲线会呈现出阶梯状的「突变」。
但如果改用\textbf{连续}的评估指标
（如 token-level 的对数似然、Brier 分数），
同一个模型的能力增长曲线就变得平滑了：
不存在明显的「突变点」。

这就好比考试成绩从 59 分跳到 60 分，
在「是否及格」这个离散指标上是一个突变，
但在连续分数上只是 1 分的平滑进步。

\begin{warnbox}[涌现之争的现状]
对涌现现象的正确理解是：
\begin{itemize}[noitemsep]
  \item 在\textbf{离散指标}（如 exact-match accuracy）下，涌现的阶跃特征确实存在
  \item 在\textbf{连续指标}（如 perplexity）下，能力增长通常是平滑的
  \item 但\textbf{实际应用}往往关心的就是离散指标（答案对不对），
        所以涌现在应用层面是「真实」的
  \item 大模型能做到小模型做不到的事，这一事实无可争议，
        争议仅在于「过渡」是突变还是平滑的
\end{itemize}
\end{warnbox}

从工程角度来说，无论涌现的「阶跃」是真实的还是度量假象，
规模仍然是最可靠的能力提升路径：
尽管 Scaling Laws 的回报率可能在递减。

\subsection{涌现的可能机制}

涌现背后的机制是什么？虽然我们还没有完整的理论解释，但几个方向的研究提供了部分洞察。

\textbf{电路假说}。Olsson 等人~\cite{olsson2022induction} 发现，Transformer 在训练过程中会形成特定的「电路」（circuit），多个注意力头协作完成特定的子任务。例如，一种被称为「induction head」的电路模式由两个注意力头组成：第一个头识别当前 token 之前出现过的模式，第二个头复制该模式后面跟着的 token。这种电路在训练过程中会突然形成，在形成之前，模型无法做模式匹配；形成之后，模式匹配能力突然出现。这种电路的形成可能是涌现的微观机制之一。

\textbf{能力的组合爆炸}。另一种假说是：复杂能力是简单子能力的组合。一个 10B 模型可能已经学会了「理解数字含义」和「执行单步加法」，但还没有学会将两者串联起来做多步算术。当模型规模增加到 100B 时，它可能学会了「将多个子技能串联」这一元能力，这导致了多步推理的突然涌现。从这个角度看，涌现不是一个全新能力的产生，而是已有子能力的\textbf{组合}达到了某个临界点~\cite{arora2023theory}。

\textbf{信息整合的门槛}。最近的研究~\cite{wei2023larger} 提出了一个更具体的假说：模型需要达到一定的规模才能在其内部表示中\textbf{同时}维持多个相关的概念。一个小模型可能能分别理解「巴黎是法国首都」和「法国是欧洲国家」，但无法在同一个推理步骤中同时激活这两个知识来回答「巴黎在哪个大洲」。更大的模型有更宽的隐藏维度和更多的注意力头，可以同时维持更多的信息，当这个容量超过某个任务的需求时，该任务的能力就「涌现」了。

这三种解释并不互斥，它们可能在不同的任务和规模范围内都有贡献。对涌现机制的理解不仅是一个科学问题，还有直接的工程价值：如果我们知道什么条件触发涌现，就可以更有针对性地设计训练策略，而不是盲目地增大规模。

\section{Scaling Laws：精确的幂律}
\label{sec:scaling-laws}

涌现现象的存在自然引出一个问题：模型能力与规模之间的关系能否被精确量化？Scaling Laws 正是对这个问题的回答。

\subsection{Kaplan 的先驱工作}

2020 年，OpenAI 的 Kaplan 等人~\cite{kaplan2020scaling} 发表了一项开创性研究，首次揭示了语言模型性能与三个核心变量之间存在惊人地平滑的\textbf{幂律关系}（power law）：

\begin{equation}
  L(N) \propto N^{-\alpha_N}, \quad L(D) \propto D^{-\alpha_D}, \quad L(C) \propto C^{-\alpha_C}
\end{equation}

其中 $L$ 是交叉熵损失，$N$ 是模型参数量，$D$ 是训练数据量（token 数），$C$ 是训练计算量（FLOPs）。指数 $\alpha$ 决定了 scaling 的效率。

Kaplan 团队的关键发现包括：

\textbf{参数量主导论}。在固定计算预算下，Kaplan 认为最优策略是尽可能增大模型参数量，同时使用相对较少的训练数据。他们给出的最优比例关系是：$N \propto C^{0.73}$，$D \propto C^{0.27}$。换句话说，当计算预算增加 10 倍时，应该将模型增大约 5.4 倍，但只增加约 1.9 倍的训练数据。

\textbf{幂律关系横跨多个数量级}。从 $10^3$ 到 $10^{10}$ 参数的模型，损失与参数量的对数-对数图几乎是一条直线。这种跨越七个数量级的规律性在自然科学中很常见（如开普勒定律），但在机器学习中极为罕见。

\textbf{架构细节相对不重要}。在同等参数量下，模型的宽度（hidden size）和深度（层数）的具体比例对性能的影响远小于参数总量本身。这是一个令人意外的发现，大量的架构搜索工作可能价值有限。

\subsection{Chinchilla 的颠覆：数据同样重要}

2022 年，DeepMind 的 Hoffmann 等人~\cite{hoffmann2022chinchilla} 发表了 Chinchilla 论文，对 Kaplan 的结论做出了重要修正。这篇论文的影响可以说改变了整个行业的训练策略。

Hoffmann 团队的方法论与 Kaplan 不同。Kaplan 团队在各种规模上训练了固定 epoch 数的模型；而 Hoffmann 团队为每个计算预算找到使最终损失最小化的参数量和数据量组合。这看似微小的方法论差异导致了截然不同的结论。

Chinchilla 定律的核心发现是：\textbf{模型参数量和训练数据量应该等比例增长}。最优比例约为每个参数对应 20 个训练 token。用公式表示：

\begin{equation}
  N_{\text{opt}} \propto C^{0.5}, \quad D_{\text{opt}} \propto C^{0.5}
\end{equation}

这与 Kaplan 的结论有质的不同。Kaplan 主张「把钱花在更大的模型上」，Chinchilla 主张「模型和数据同等重要」。

\textbf{为什么 Kaplan 错了？}差异的根源在于实验设计。Kaplan 团队使用固定的学习率调度器（对所有规模使用相同的训练步数），而 Chinchilla 团队为每个计算预算独立优化了训练步数和学习率。Kaplan 的实验中，大模型因为训练步数相对不足而「看起来」效率更高，如果给小模型同样充足的训练，差距会大幅缩小~\cite{hoffmann2022chinchilla}。

\textbf{实际影响}是巨大的。按照 Kaplan 的建议，GPT-3（175B）的训练数据量（300B tokens）是合理的。但按照 Chinchilla 定律，175B 的模型应该用约 3.5T tokens 来训练——GPT-3 的训练数据\textbf{远远不足}。Chinchilla（70B 参数，1.4T tokens）正是这一发现的直接产物：它的参数量只有 Gopher（280B）的四分之一，但因为用了充足的训练数据，性能反而更强。

这个发现对整个行业的影响堪称地震级别。在 Chinchilla 论文之前，OpenAI、Google 和各大实验室的主要策略是「把模型做大」，把计算预算花在更多的参数上。Chinchilla 之后，策略变为「在参数和数据之间找到最优平衡」。Meta 的 LLaMA~\cite{touvron2023llama} 是这一转变的标志性产物：LLaMA-7B 用 1T tokens 训练，LLaMA-13B 也用 1T tokens 训练——远超 Chinchilla 定律建议的数据量，目的是让小模型也能达到竞争力强的性能。这种「过度训练」（over-training）策略在 2023--2024 年成为主流：既然推理成本与参数量线性相关，那么训练一个小但充分训练的模型，在总拥有成本上可能比训练一个大但训练不足的模型更经济。

\subsection{Chinchilla 之后：规模定律的持续演化}

Chinchilla 定律虽然影响巨大，但它只是故事的开始。后续研究揭示了更多维度的 scaling 行为。

\textbf{推理时计算的 Scaling}。Chinchilla 定律只考虑了训练计算量，完全忽略了推理成本。但在实际部署中，推理成本往往远超训练成本，一个模型训练一次但被调用数十亿次。Sardana 和 Frankle~\cite{sardana2024chinchilla} 提出了「推理最优」的 Chinchilla 定律：如果考虑模型的总生命周期成本（训练 + 推理），最优的模型应该更小但训练更久。这解释了为什么 LLaMA~\cite{touvron2023llama} 选择用 1.4T tokens 训练一个 7B 模型——远超 Chinchilla 的「最优」比例，但在推理成本上更经济。

\textbf{数据质量的 Scaling}。2024--2025 年的研究表明，数据质量的影响可能与数据数量同等重要。在相同的 token 数下，经过精心清洗和去重的数据集可以带来等价于 2--5 倍数据量增加的性能提升~\cite{penedo2024fineweb}。这意味着 Chinchilla 定律中的 $D$ 不应该被简单地理解为 token 数，而应该是某种「有效数据量」——考虑了数据质量、多样性和去重程度。

\textbf{后训练的 Scaling}。传统 Scaling Laws 只关注预训练损失。但 DeepSeek-R1~\cite{deepseekr1} 的成功表明，后训练（特别是强化学习）阶段也存在独立的 scaling 行为，更多的 RL 计算可以持续提升推理能力，即使预训练模型相同。这是一个全新的 scaling 维度，传统理论尚未完全覆盖。

\subsection{Scaling Laws 的统一视角}

将 Kaplan、Chinchilla 以及后续的研究放在一起，一个更完整的图景浮现了。模型的最终能力取决于一个\textbf{多维 scaling 空间}中的资源分配：

\begin{equation}
  \text{Performance} = f(N, D, Q, C_{\text{train}}, C_{\text{post}}, C_{\text{test}})
\end{equation}

其中 $N$ 是参数量，$D$ 是数据 token 数，$Q$ 是数据质量，$C_{\text{train}}$ 是预训练计算量，$C_{\text{post}}$ 是后训练计算量，$C_{\text{test}}$ 是推理时计算量。传统的 Scaling Laws 只考虑了前四个变量；Chinchilla 引入了 $N$ 和 $D$ 之间的最优平衡；后训练 scaling 和 test-time compute 又打开了两个新维度。

这个多维视角解释了为什么不同的团队在相近的性能水平上做出了截然不同的设计选择：
\begin{itemize}[noitemsep]
  \item \textbf{DeepSeek-V3}：中等 $N$（671B 总/37B 活跃），大 $D$（14.8T tokens），高 $Q$（精细的数据工程），中等 $C_{\text{post}}$
  \item \textbf{GPT-5}：大 $N$（未公开，推测 $>$ 1T），大 $D$，大 $C_{\text{post}}$（推理路由），大 $C_{\text{test}}$（自动深度思考）
  \item \textbf{DeepSeek-R1}：中等 $N$（同 V3 基座），中等 $D$，极大 $C_{\text{post}}$（纯 RL 训练），大 $C_{\text{test}}$（长思维链）
\end{itemize}

每种组合都代表了在多维 scaling 空间中的不同位置，而 frontier 正是这些不同位置的包络线。

\begin{warnbox}[Scaling Laws 的局限性]
Scaling Laws 是经验性的幂律拟合，不是物理定律。它们有几个重要的局限：
\begin{itemize}[noitemsep]
  \item \textbf{外推风险}：幂律在已验证的范围内成立，但没有理论保证它们在更大规模上仍然成立。可能存在未知的拐点或饱和效应。
  \item \textbf{损失 $\neq$ 能力}：Scaling Laws 预测的是交叉熵损失，不是下游任务的表现。损失的平滑下降并不排除特定能力的非线性涌现。
  \item \textbf{数据分布依赖}：幂律的具体参数（$\alpha$ 值）强烈依赖于训练数据的分布。在代码、数学等特定领域，scaling 行为可能显著不同。
  \item \textbf{架构依赖}：Scaling Laws 是在特定架构家族（标准 Transformer）上测量的。MoE、SSM（如 Mamba）等新架构可能有不同的 scaling 行为，但获取这些数据需要大量的实验，目前的证据仍然有限。
\end{itemize}
\end{warnbox}

\subsection{Scaling Laws 作为工程决策工具}

尽管有上述局限，Scaling Laws 在实践中仍然是训练大模型最重要的规划工具。以下是它的几种典型应用场景。

\textbf{训练前的成本估算}。在启动一次数百万美元的训练之前，团队需要回答一个关键问题：「在给定的计算预算下，我们能期望什么水平的模型性能？」Scaling Laws 提供了一个量化的答案。例如，如果目标是 GPT-4 级别的 perplexity，可以用已知的 scaling 曲线反推需要的参数量和训练数据量——进而估算 GPU 小时和成本。

\textbf{最优模型配置}。给定固定的计算预算 $C$，如何选择模型大小 $N$ 和训练数据量 $D$？Chinchilla 定律给出了一个精确的答案：$N \approx \sqrt{C/(6 \cdot 20)}$，$D = 20N$。但这是「训练最优」而非「部署最优」，如果考虑推理成本，最优的 $N$ 会更小（因为小模型推理更便宜），相应地 $D$ 会更大（用更多数据来弥补参数的减少）。

\textbf{小规模实验的外推}。Scaling Laws 最实用的应用之一是：在小规模上做实验，然后外推到大规模。例如，如果一种新的数据清洗方法在 1B 模型上将 perplexity 降低了 5\%，Scaling Laws 可以帮助预测这种改进在 70B 模型上是否仍然有效、改进幅度如何。这种「小到大」的外推虽然不是完全可靠的，但它极大地减少了大规模实验的试错成本——一次小规模实验只需要几千美元，而一次大规模实验可能需要几百万美元。

\textbf{调试训练过程}。训练过程中，如果损失曲线的下降速度偏离了 Scaling Laws 的预测，这通常意味着出了问题，可能是数据管道的 bug、学习率设置不当、或者硬件故障导致的通信错误。Scaling Laws 在这里充当了训练过程的「基准线」，帮助工程师快速定位问题。

\section{LLM 训练的完整流水线}

在深入各个技术细节之前，让我们先建立一个全局视角。从原始文本到可部署的 AI 系统，LLM 训练是一条多阶段的流水线。每个阶段都有独特的技术挑战和设计考量。

\begin{figure}[htbp]
\centering
\begin{tikzpicture}[
  node distance=10mm and 18mm,
  stage/.style={
    rectangle, rounded corners=3pt,
    draw=#1, line width=0.9pt,
    fill=#1!6, minimum width=26mm, minimum height=14mm,
    inner sep=6pt, font=\small\bfseries,
    text=black!80, align=center
  },
  sub/.style={
    font=\scriptsize, text=figGray, align=center, inner sep=0pt
  },
  arr/.style={-{Stealth[length=3pt,width=3pt]}, line width=0.7pt, figGray!70},
  phase/.style={
    rounded corners=5pt, draw=none, fill=#1!6, inner xsep=8pt, inner ysep=14pt
  },
  plabel/.style={
    font=\scriptsize\sffamily\bfseries, text=#1, fill=white,
    inner sep=2pt, rounded corners=1.5pt
  },
]

% === Phase backgrounds (drawn first, behind everything) ===
\begin{scope}[on background layer]
  \fill[figBlue!5, rounded corners=5pt]
    (-8mm, -12mm) rectangle (148mm, 20mm);
  \fill[figOrange!5, rounded corners=5pt]
    (-8mm, -58mm) rectangle (148mm, -24mm);
\end{scope}

% === Top row: Data Engineering + Pre-training ===
\node[stage=figBlue] (data) at (0,0) {数据采集};
\node[sub, below=2pt of data] {Common Crawl, 书籍, 代码};

\node[stage=figBlue, right=18mm of data] (clean) {清洗与去重};
\node[sub, below=2pt of clean] {质量评分, MinHash 去重};

\node[stage=figBlue, right=18mm of clean] (tok) {分词};
\node[sub, below=2pt of tok] {BPE/Unigram, 32K--262K};

\node[stage=figGreen!80!black, right=18mm of tok] (pre) {预训练};
\node[sub, below=2pt of pre] {Next-token, $10^{13}$ tokens};

% === Bottom row: Post-training + Deployment ===
\node[stage=figOrange, below=36mm of data] (sft) {SFT};
\node[sub, below=2pt of sft] {指令跟随, 对话};

\node[stage=figOrange, right=18mm of sft] (rl) {强化学习};
\node[sub, below=2pt of rl] {RLHF / GRPO / DPO};

\node[stage=figRed!80!black, right=18mm of rl] (eval) {评估};
\node[sub, below=2pt of eval] {基准测试, 人类评估};

\node[stage=figRed!80!black, right=18mm of eval] (deploy) {部署};
\node[sub, below=2pt of deploy] {量化, vLLM, 蒸馏};

% === Arrows ===
\draw[arr] (data) -- (clean);
\draw[arr] (clean) -- (tok);
\draw[arr] (tok) -- (pre);
\draw[arr, rounded corners=4pt] (pre.south) -- ++(0,-10mm) -| (sft.north);
\draw[arr] (sft) -- (rl);
\draw[arr] (rl) -- (eval);
\draw[arr] (eval) -- (deploy);

% === Phase labels ===
\node[plabel=figBlue, anchor=north west] at (-4mm, 16mm)
  {\small 数据密集型阶段};
\node[plabel=figOrange, anchor=north west] at (-4mm, -28mm)
  {\small 精细化阶段};

% === Feedback arrow (eval -> SFT) ===
\draw[arr, dashed, figGray!50, rounded corners=3pt]
  (deploy.south) -- ++(0,-4mm) -| ([xshift=-4mm]sft.south)
  node[pos=0.5, below, sub] {迭代};

\end{tikzpicture}
\caption{LLM 训练流水线全景。上排为数据密集型阶段（数据采集至预训练），
下排为精细化阶段（SFT 至部署）。垂直箭头标志着从"大规模无监督"到"小规模有监督"的范式转换。}
\label{fig:training-pipeline}
\end{figure}

\subsection{数据工程阶段}

整条流水线的第一步是数据。没有数据，一切都是空谈。

\textbf{数据采集}是一项被严重低估的工程挑战。主要数据来源包括 Common Crawl（互联网爬取数据，约 250B 页面）、Wikipedia（高质量百科知识）、GitHub（代码数据）、学术论文（arXiv、Semantic Scholar）以及书籍语料。不同来源的数据质量差异巨大——Common Crawl 中超过 90\% 的内容是低质量的：广告、垃圾页面、重复内容、机器生成的 SEO 文本。

\textbf{清洗与去重}是数据工程中最关键的步骤。典型的清洗流水线包括：语言识别（过滤非目标语言的文档）、质量评分（用小型分类器评估文档质量）、内容过滤（移除有害内容、个人信息）、模糊去重（使用 MinHash/SimHash 移除近似重复的文档）。FineWeb~\cite{penedo2024fineweb} 的经验表明，清洗质量的提升可以带来等价于数倍数据量增加的性能改善——这是数据工程被低估的一个证据。

去重的重要性怎么强调都不过分。Lee 等人的研究表明，在未去重的 Common Crawl 上训练的模型，其有效训练数据量可能只有标称数据量的 30--50\%，因为大量文档是近似重复的（模板化网页、镜像站点、爬虫陷阱）。重复数据不仅浪费训练计算，还会导致模型\textbf{过度拟合}重复出现的模式，降低泛化能力。更微妙的是，重复数据中的错误会被模型以更高的权重学习，如果某个错误的事实陈述出现了 100 次，模型学到的是「这个说法很可靠」，而非「这个说法有争议」。

\textbf{数据配比}（data mixing）是一个越来越被重视的设计决策。训练数据通常来自多个来源——网页、书籍、代码、学术论文、对话数据，它们的质量和主题分布各不相同。如何设置各来源的采样比例直接影响模型的能力分布。例如，增加代码数据的比例会提升代码生成能力，但可能以自然语言理解的轻微下降为代价。DeepSeek-V3~\cite{deepseekv3} 在训练的不同阶段使用了不同的数据配比——早期以网页数据为主建立通用知识，后期增加代码和学术数据提升专业能力。这种\textbf{课程学习}（curriculum learning）策略是 2025--2026 年主流训练方案的标准做法。

\textbf{分词}将清洗后的文本转换为模型可处理的 token 序列。这个看似简单的步骤对模型的多语言能力、数学能力和代码能力都有深远影响，我们将在第~\ref{ch:tokenization}章详细讨论。

\subsection{预训练阶段}

预训练是整条流水线中计算成本最高的阶段，也是模型核心能力的来源。

在这个阶段，模型在数万亿 token 上进行 next-token prediction 训练。这需要数千块 GPU 运行数周到数月。DeepSeek-V3 使用 2048 块 H800 GPU 训练了约 2 个月~\cite{deepseekv3}；GPT-4 的训练据推测动用了超过 25,000 块 A100 GPU。

预训练的技术挑战涵盖了模型架构设计（第~\ref{ch:transformer}章）、分布式训练策略（第~\ref{ch:pretrain}章）、混合精度训练、学习率调度以及训练稳定性等诸多方面。一次训练失败的成本可能高达数百万美元，这使得训练稳定性成为一个工程上的关键问题。

\textbf{分布式训练}是预训练的工程核心。一个 175B 参数的模型以 FP16 存储需要 350GB 内存，远超单张 GPU 的容量（A100 最大 80GB）。解决方案是将模型和数据分布到数千块 GPU 上。主流的并行策略包括：
\begin{itemize}[noitemsep]
  \item \textbf{数据并行}（Data Parallelism）：每块 GPU 保存模型的完整副本，处理不同的数据批次，然后同步梯度。ZeRO 优化~\cite{rajbhandari2020zero} 将优化器状态、梯度和模型参数分片存储，大幅降低了内存需求。
  \item \textbf{张量并行}（Tensor Parallelism）：将模型的每一层切分到多块 GPU 上，每块 GPU 负责层内计算的一部分。这种方式需要 GPU 之间的高带宽通信（通常是 NVLink 或 NVSwitch）。
  \item \textbf{流水线并行}（Pipeline Parallelism）：将模型的不同层分配到不同的 GPU 上，数据像流水线一样从前向后流过。这减少了对 GPU 间带宽的需求，但引入了「气泡」（idle time）问题。
  \item \textbf{序列并行}（Sequence Parallelism）：将长序列切分到多块 GPU 上，每块 GPU 处理序列的一部分。这在处理超长上下文时尤为重要。
\end{itemize}

实际的训练方案通常是这四种并行策略的\textbf{混合}，例如，DeepSeek-V3 同时使用了数据并行、张量并行和流水线并行，配合 FP8 混合精度训练，在 2048 块 H800 GPU 上实现了极高的计算效率。

\textbf{训练稳定性}是另一个关键挑战。在数千块 GPU 上训练数周，任何一次梯度爆炸（loss spike）、硬件故障或通信错误都可能导致训练中断。主流的稳定性技巧包括：梯度裁剪（gradient clipping）、学习率预热（warmup）和衰减（decay）、损失尖峰自动检测和回滚（rollback）、以及定期保存检查点（checkpoint）。DeepSeek-V3 在整个训练过程中没有遇到需要回滚的损失尖峰，这被认为是其训练工程质量的一个重要标志。

训练稳定性的另一个维度是\textbf{可复现性}。
如果在完全相同的硬件和数据上重新训练，
模型是否能达到相同的性能？
由于浮点运算的非结合性（$a + (b + c) \neq (a + b) + c$）
和分布式训练中 all-reduce 操作的不确定性，
大模型训练通常不是比特级可复现的：
但在统计意义上应该是可复现的
（损失曲线在 $\pm 1\%$ 的范围内重合）。
如果两次训练的结果显著不同，
通常意味着存在隐藏的 bug
（如数据加载顺序的非确定性、
随机数种子的不一致等）。

\textbf{混合精度训练}是现代预训练的标配。
FP32（32 位浮点）太慢也太耗内存；
FP16（16 位浮点）足够快，但数值范围有限，
容易导致梯度下溢（underflow）。
BF16（Brain Float 16）保持了 FP32 的指数范围
但减少了尾数精度——这是一个极好的折中。
DeepSeek-V3 更进一步，使用了 FP8 训练：
将精度进一步压缩到 8 位。
这需要精心的混合精度策略：
前向传播和反向传播在 FP8 中进行，
而优化器状态（动量和方差）保持在 FP32：
因为优化器状态需要高精度来累积小梯度。
FP8 训练将内存占用减半、计算吞吐量翻倍，
是 DeepSeek-V3 能够以 560 万美元达到 GPT-4 性能的关键技术之一。

这些技术细节将在第~\ref{ch:pretrain}章中展开讨论。
这里重要的是建立一个直觉：
\textbf{预训练不仅仅是「把数据喂给模型」：
它是一项涉及硬件、通信、数值精度、容错、调度等多个维度的系统工程}。
一个看似微小的工程决策
（如选择 BF16 而非 FP16，或选择 ZeRO-3 而非张量并行）
可能导致训练效率的数倍差异。
这也是为什么「工程能力」在 LLM 竞赛中与「算法创新」同等重要。
DeepSeek 团队在技术报告中自豪地声称其训练过程零次回滚：
这背后是数月的基础设施投入和严格的工程实践。

\subsection{后训练阶段}

预训练产出的是一个「基座模型」（base model），它具备强大的语言能力，但行为不可控，可能生成有害内容，也不擅长遵循用户指令。后训练的目标是将基座模型转变为一个有用、安全、对齐的助手。

\textbf{监督微调}（SFT）使用人类标注的指令-响应对来教模型遵循指令。高质量的 SFT 数据集通常只需要数万到数十万条样本——与预训练的万亿 token 相比，数据量小了好几个数量级，但对模型行为的影响却是决定性的。

\textbf{强化学习}（RLHF 或其他变体）使用人类偏好信号来进一步优化模型。传统的 RLHF~\cite{ouyang2022rlhf} 先训练一个奖励模型，再用 PPO 算法优化策略。DeepSeek-R1~\cite{deepseekr1} 证明了更简洁的 RL 方案（如 GRPO）同样有效，甚至可以仅凭 RL 就让模型学会思维链推理。

\subsection{后训练阶段的关键转折}

后训练阶段在 2024--2025 年经历了一次方法论的根本转变。传统的 RLHF 流水线（训练奖励模型 $\to$ PPO 优化）被证明并非唯一选择，甚至不一定是最优选择。

DeepSeek-R1~\cite{deepseekr1} 展示了一条令人意外的路径：直接使用 GRPO（Group Relative Policy Optimization）进行强化学习，不需要单独的奖励模型，也不需要人类标注的推理过程。模型在 RL 训练过程中\textbf{自发地}学会了思维链推理，包括反思、回溯和自我纠正。这个发现的意义在于：高级推理能力不需要通过模仿人类来学习，它可以从简单的结果反馈（答案对不对）中自然涌现。

DPO（Direct Preference Optimization）~\cite{rafailov2023dpo} 代表了另一个简化方向：完全消除 RL 的在线采样环节，将偏好学习重新表述为一个简单的二元分类问题。DPO 在工程实现上比 PPO 简单得多（不需要维护四个模型的副本），训练也更稳定。但 DPO 是否能匹配 PPO 在最难任务上的效果，仍是活跃的研究话题。

\subsection{评估与部署}

训练完成的模型需要经过严格的评估才能部署。评估涵盖通用能力（MMLU、HellaSwag）、推理能力（GSM8K、MATH）、代码能力（HumanEval、SWE-bench）、安全性（TruthfulQA、SafetyBench）等多个维度。

2025--2026 年，评估本身也在经历变革。传统的静态基准（MMLU、GSM8K）在 frontier 模型上已经趋于饱和——GPT-5 和 Claude Opus 4.6 在 MMLU 上的分数超过 90\%，接近人类专家水平。新一代评估基准开始转向更具挑战性的方向：Humanity's Last Exam 汇集了各领域专家提出的「最难的问题」；SWE-bench Verified 测试模型解决真实 GitHub issue 的能力；FrontierMath~\cite{glazer2024frontiermath} 包含了研究级数学难题。这些新基准的共同特点是：它们测试的不是记忆，而是\textbf{真正的推理和问题解决能力}。

部署阶段的核心挑战是推理效率。一个 175B 参数的模型在全精度下需要约 350GB 的 GPU 内存，这超出了任何单张 GPU 的容量。量化（将权重从 FP16 压缩到 INT8 或 INT4）、蒸馏（用大模型训练小模型）、投机解码（speculative decoding）等技术使得高效部署成为可能。

MoE 架构在部署中展现了独特的优势：DeepSeek-V3 的 671B 总参数中每次推理只激活 37B，这意味着它在推理成本上接近一个 37B 的密集模型，但能力却接近 671B。这种「以空间换时间」的策略是 2025--2026 年主流部署方案的核心。

\subsection{各阶段的成本分布}

理解 LLM 训练流水线各阶段的相对成本，有助于合理分配资源。以一次 frontier 模型训练为例（如 DeepSeek-V3 级别），各阶段的成本占比大致如下：

\begin{center}\small
\renewcommand{\arraystretch}{1.3}
\begin{tabular}{lrl}
\toprule
\rowcolor{figLightGray}
\textbf{阶段} & \textbf{成本占比} & \textbf{时间尺度} \\
\midrule
数据采集与清洗    & 5--10\%   & 数周到数月 \\
分词器训练        & $<$1\%    & 数小时 \\
预训练            & 70--85\%  & 数周到数月 \\
SFT               & 2--5\%    & 数天 \\
RL（RLHF/GRPO）   & 5--15\%   & 数天到数周 \\
评估              & 1--3\%    & 数天 \\
\bottomrule
\end{tabular}
\end{center}

这个分布揭示了一个关键事实：\textbf{预训练占据了绝大部分成本，但后训练对最终用户体验的影响可能更大}。一个未经后训练的 base model 虽然在语言建模指标上可能很强，但在实际使用中几乎不可用，它不会遵循指令，不会拒绝有害请求，生成的文本风格也不符合人类偏好。这就是为什么后训练虽然只占总成本的 10--20\%，却值得投入大量的研究和工程精力。

DeepSeek-R1~\cite{deepseekr1} 的训练方式更极端地体现了这一点：它的预训练阶段直接复用了 V3 的基座模型（零额外成本），全部的创新都在后训练阶段，通过纯 RL 训练让模型学会推理。这意味着 R1 的「边际成本」几乎全部在 RL 阶段，而这个阶段产出的推理能力恰恰是 R1 最大的卖点。

\section{从 GPT 到 DeepSeek：一部效率进化史}

LLM 的历史可以粗略地分为几个阶段，
但在进入时间线之前，
让我们先理解一个贯穿整个历史的核心转变：
从\textbf{特征工程}到\textbf{上下文学习}，
从\textbf{任务专用}到\textbf{通用基础模型}。

\subsection{前传：词向量与上下文表示（2013--2018）}

在 Transformer 之前，自然语言处理经历了两次重要的范式转变。

\textbf{Word2Vec}~\cite{mikolov2013word2vec}（2013）和 GloVe（2014）
证明了一个惊人的事实：
通过在大规模语料上训练词向量，
语义关系可以被编码为向量空间中的方向。
但这些是\textbf{静态}词向量：
无论上下文如何，「bank」这个词只有一个表示，
无法区分「银行」和「河岸」。

\textbf{ELMo}~\cite{peters2018elmo}（2018）解决了这个问题：
通过双向 LSTM 生成\textbf{上下文相关}的词表示。
同一个词在不同句子中会有不同的向量。
ELMo 的出现标志着 NLP 领域从静态表示走向动态表示，
为后来的预训练范式铺平了道路。

\subsection{架构探索期（2017--2019）}

2017 年，Transformer~\cite{vaswani2017attention} 被提出，
最初只是一个机器翻译模型。
但它的自注意力机制带来的并行化优势
很快引起了整个领域的注意。

Transformer 的出现是一个历史性的转折点，
但它的重要性并非立刻被认识到。
论文标题「Attention Is All You Need」在当时被许多研究者视为夸大其词：
注意力机制早在 2014 年就已存在（Bahdanau attention），
Transformer 只是将其推到了极致。
然而事后看来，Transformer 的关键创新并非注意力机制本身，
而是\textbf{完全基于注意力的架构}：
去掉了 RNN 的顺序依赖，实现了完全的并行化。
这一设计选择使得 Transformer 可以充分利用 GPU 的并行计算能力：
当模型需要在数千块 GPU 上训练时，这种并行性是生死攸关的。

一个常被忽略的事实是：
Transformer 论文的第一作者是 Ashish Vaswani，
而非 AI 领域更知名的研究者。
论文的八位作者中有六位后来离开了 Google：
创办了 Character.AI、Cohere、Inceptive 等公司，
或加入了 OpenAI。
这个「transformer 八人组」的故事本身就说明了一个问题：
\textbf{最重要的研究突破并不总是来自最资深的研究者，
但一旦突破发生，其影响会远超原始论文的范畴}。

2018 年发生了两件大事：

\textbf{GPT-1}（2018, 117M 参数）选择了\textbf{decoder-only}路线：
单向语言模型，通过预测下一个 token 来预训练，
再在下游任务上微调。
它首次证明了「预训练 + 微调」范式的威力。

\textbf{BERT-Large}~\cite{devlin2019bert}（2018, 110M/340M 参数）选择了\textbf{encoder-only}路线：
双向语言模型，通过 Masked Language Model（随机遮盖 15\% 的 token 并预测）来预训练。
BERT 在 11 项 NLP 基准上刷新了记录，
震动了整个学界。

这一阶段的模型较小，训练成本以千美元计。
两条路线的竞争刚刚开始。

同时期还有一条 encoder-decoder 路线：
\textbf{T5}~的前身（2018--2019）探索了
「将所有 NLP 任务统一为 text-to-text」的框架。
这三条路线——encoder-only（BERT）、decoder-only（GPT）、
encoder-decoder（T5/BART），在 2019--2020 年展开了激烈竞争。
竞争的结果到 2022 年变得清晰：
对于文本生成和通用任务，decoder-only 路线全面胜出。
BERT 虽然在特定的理解任务上仍有价值，
但无法扩展到生成场景；
encoder-decoder 在翻译和摘要上有优势，
但在工程简洁性上输给了 decoder-only。

\textbf{GPT-2}（2019, 1.5B 参数）是一个转折点。
OpenAI 声称它「太危险而不能公开发布」：
这在当时引发了激烈争议，但也将公众注意力吸引到了语言模型上。
更重要的是，GPT-2 展示了零样本能力的雏形：
不需要微调，仅通过提示就能完成多种任务。
回顾历史，GPT-2 的「危险」论并非完全没有道理：
它确实是第一个能生成通顺、连贯、有时令人信服的长文本的模型。
但更深远的影响是它预示了一种新的 AI 使用范式：
不再需要为每个任务收集数据、训练模型、部署推理：
一个通用的语言模型就能「做一切」。
这个愿景在 GPT-3 中初步实现，在 GPT-4 中基本实现，
到了 GPT-5 时代已经成为日常现实。

\subsection{规模竞赛期（2020--2022）}

\textbf{GPT-3}~\cite{brown2020gpt3}（2020, 175B 参数）
是一个划时代的模型。
它的训练成本约 460 万美元，参数量比 GPT-2 大 100 倍。
GPT-3 最重要的贡献是发现了\textbf{in-context learning}（上下文学习）：
你不需要微调模型，只需在提示中给出几个示例，
模型就能学会新任务。
这从根本上改变了 NLP 的使用范式：
从「训练一个任务专用模型」变成「设计好的提示给通用模型」。

In-context learning 的发现是偶然的还是必然的？
回顾历史，GPT-3 的研究者并非有意设计了这种能力：
它是规模增大的\textbf{副产品}。
GPT-2（1.5B）已经展示了零样本能力的微弱迹象，
但效果太不稳定，研究者没有深入追踪。
GPT-3 将参数量增大了 100 倍后，
in-context learning 变得足够可靠，以至于可以被系统地研究和利用。
这个故事完美地说明了涌现的本质：
\textbf{某些能力只有在足够大的规模上才变得显著，
在此之前它们隐藏在噪声中}。

GPT-3 还引发了一场关于「理解」的哲学争论。
批评者指出，GPT-3 只是在做「随机鹦鹉」式的模式匹配，
并不真正「理解」语言。
支持者反驳说，如果模式匹配能做到如此之多：
翻译、推理、写诗、编程：
那么「真正的理解」和「足够好的模式匹配」之间的区别是什么？
这场辩论至今仍在进行，
但一个实用主义的共识正在形成：
无论模型是否「真正理解」，
只要它能可靠地完成人类关心的任务，
关于「理解」的哲学问题就变得次要了。

从工程的角度看，这场辩论的实际意义在于：
它影响了我们对模型能力边界的预期。
如果 LLM 只是在做模式匹配，
那么它在分布外的任务上应该表现糟糕。
如果它在某种程度上学会了世界模型，
那么它应该能处理全新的情境。
GPT-4 和后续模型在各种新颖的推理任务上的表现
越来越倾向于后一种解释：
但确凿的证据仍然不足。
这个问题的最终答案可能需要
更深入的理论工作和更精心设计的实验。

Google 在这一时期也没有闲着。
\textbf{T5}（2020, 11B 参数）提出了
「将所有 NLP 任务统一为 text-to-text」的框架：
无论是分类、翻译还是摘要，都是输入一段文本、输出一段文本。
\textbf{PaLM}~\cite{chowdhery2022palm}（2022, 540B 参数）
是当时最大的密集模型，展示了规模带来的能力跃升，
特别是在推理和多语言任务上。

DeepMind 的 \textbf{Chinchilla}~\cite{hoffmann2022chinchilla}（2022, 70B 参数）
带来了一个关键发现：之前的模型普遍\textbf{训练不足}。
Chinchilla 定律指出，模型参数和训练数据应该等比例增长：
一个 70B 的模型如果用足够的数据训练（1.4T tokens），
可以超越参数量大 4 倍但数据不足的 280B 模型。
这深刻地影响了后续所有模型的训练策略。

训练成本在这一时期上升到数千万美元级别。

\subsection{产品化爆发期（2023--2024）}

2022 年 11 月，\textbf{ChatGPT} 的发布引发了全球性的 AI 浪潮。
ChatGPT 本质上是 GPT-3.5 经过 RLHF（基于人类反馈的强化学习）
~\cite{ouyang2022rlhf}微调的产物：
技术上并非革命性突破，
但\textbf{对话式交互}的产品形态让普通人第一次直观感受到了 AI 的能力。
它在两个月内获得了 1 亿用户，
创造了消费者应用增长最快的历史记录。

\textbf{GPT-4}（2023）将能力推向了新高度。
虽然 OpenAI 没有公开其架构细节，
但广泛推测它采用了 MoE（混合专家）架构，
总参数量可能超过 1T，但每次推理只激活其中一部分。
GPT-4 在律师资格考试中排名前 10\%，
能够理解图片，生成复杂代码。

与此同时，开源运动蓬勃兴起。
Meta 在 2023 年发布 \textbf{LLaMA}~\cite{touvron2023llama}（7B--65B），
展示了较小模型经过充分训练后也能达到优秀性能。
LLaMA 的开源引爆了整个社区：
数百个微调版本在几周内涌现，
从 Alpaca 到 Vicuna 到 WizardLM。
随后的 \textbf{Llama 2}（2023）和 \textbf{Llama 3}~\cite{dubey2024llama3}（2024）
持续提升了开源模型的水平。

Google 发布了 \textbf{Gemini}~\cite{anil2023gemini}（2023），
这是一个原生多模态模型。
Anthropic 发布了 Claude 系列模型，
以安全性和长上下文能力著称。

中国的大模型生态也在这一时期快速发展。
清华的 \textbf{GLM}（ChatGLM）是最早对标 ChatGPT 的中文模型之一，
它选择了 prefix LM 架构：
一种介于 encoder-only 和 decoder-only 之间的折中方案。
阿里巴巴的 \textbf{Qwen} 系列逐步追赶国际水平，
以其对中文的深度优化（包括大词表、中文数据质量）著称。
\textbf{DeepSeek} 以其开放的技术报告和创新的架构设计
赢得了全球社区的关注：
特别是其 Multi-head Latent Attention（MLA）
被广泛认为是 2024 年最重要的架构创新之一，
后来被多个其他模型借鉴。
百度的 \textbf{文心一言}（ERNIE Bot）
和字节跳动的 \textbf{豆包}（Doubao）
则更侧重于产品化和商业落地。

2024 年 9 月，OpenAI 发布了 \textbf{o1} 预览版：
这是第一个专用推理模型，通过在回答前进行长时间内部「思考」，
在数学、编程和科学推理上取得了质的飞跃。
o1 开创了\textbf{test-time compute scaling} 的新范式，
为后续的 o3、DeepSeek-R1 以及各家模型的推理增强模式奠定了基础。

训练成本达到数亿美元级别。

\subsection{效率革命期（2025--2026）}

2024 年末至 2025 年初，效率革命的序幕由中国团队拉开。
\textbf{DeepSeek-V3}~\cite{deepseekv3}（2024 年 12 月）以 560 万美元的训练成本
达到了 GPT-4 级别的性能：
这是训练成本的\textbf{两个数量级下降}。
MoE 架构、FP8 训练、Multi-head Latent Attention（MLA）、
高效的 Multi-Token Prediction 以及极致的工程优化
共同推动了这场效率革命。

2025 年 1 月，\textbf{DeepSeek-R1}~\cite{deepseekr1} 在推理模型领域掀起了另一场革命。
它证明了通过纯强化学习（不依赖于人类标注的推理数据），
模型可以自发地学会思维链推理。
DeepSeek-R1 在数学、代码和推理任务上达到了与 OpenAI o1 相当的水平，
且完全开源（MIT 许可证）。
其蒸馏版本 R1-Distill-Qwen-32B 在多个推理基准上
甚至超越了 OpenAI o1-mini，
以极低的参数量实现了令人瞩目的推理能力。

2025 年春季迎来了一波密集发布。
4 月，OpenAI 同时推出了 \textbf{o3} 和 \textbf{o4-mini}：
第二代推理模型，首次实现了推理模型的工具调用能力，
包括网页搜索、Python 代码执行和视觉推理。
同期发布的 \textbf{GPT-4.1} 以百万 token 上下文窗口
重新定义了长上下文处理的标杆。
Meta 发布了 \textbf{Llama~4} 系列（Scout 和 Maverick），
首次采用 MoE 架构并原生支持多模态，
Scout 甚至支持 1000 万 token 的上下文窗口。
阿里巴巴发布了 \textbf{Qwen3}~\cite{qwen2025qwen3}，
首创「混合推理」模式——同一模型可在深度思考与快速响应之间无缝切换。

5 月，Anthropic 发布了 \textbf{Claude~4}（Opus~4 和 Sonnet~4），
在 agentic coding 和长时间自主任务上树立了新标杆。
8 月，OpenAI 推出了划时代的 \textbf{GPT-5}：
一个统一的系统，内置了智能路由器，
能自动判断何时快速响应、何时调用深度推理，
实质上将 GPT 系列和 o 系列合二为一。

下半年的节奏更加紧凑：
9 月，\textbf{DeepSeek-R1 登上 \textit{Nature} 封面}：
这是首篇通过独立同行评审发表在 \textit{Nature} 的 LLM 论文，
梁文锋为通讯作者，
正式确认纯强化学习可以激励模型自发产生高级推理能力。
Anthropic 的 \textbf{Claude Sonnet~4.5}（9 月）和 \textbf{Claude Opus~4.5}（11 月）
将 agentic coding 推向新高度；
Google 发布 \textbf{Gemini~3}（11 月），
以 Deep Think 模式和原生多模态能力重新加入竞争。
DeepSeek 发布 \textbf{V3.2}（12 月），引入 Sparse Attention 机制，
性能达到 GPT-5 水平，高计算变体 \textbf{V3.2-Speciale}
在 2025 年 IMO 中取得金牌（35/42 分）、ICPC 第 2 名、IOI 第 10 名。
年底，DeepSeek 连续发表两篇架构创新论文：
mHC（Manifold-Constrained Hyper-Connections，12 月 31 日）
用 Sinkhorn-Knopp 算法替代沿用十年的残差连接，
将信号放大从 3000 倍压缩至 1.6 倍；
Engram（Conditional Memory via Scalable Lookup，2026 年 1 月 12 日）
提出将「记忆」与「计算」分离的 LTM 架构：
这两项创新被普遍认为是 V4 的核心技术储备。

进入 2026 年，迭代速度进一步加快。
Anthropic 在 2 月连续发布 \textbf{Claude Opus~4.6} 和 \textbf{Sonnet~4.6}，
首次为 Opus 级模型提供百万 token 上下文窗口（beta）。
Google 推出 \textbf{Gemini~3.1~Pro}（2 月），
阿里巴巴发布 \textbf{Qwen3.5}（2 月），
以 397B 旗舰模型和全系列小模型（0.8B--9B）覆盖从云端到边缘的全场景。
OpenAI 则以 \textbf{GPT-5.3-Codex}（2 月）和 \textbf{GPT-5.4}（3 月）
持续推进 agentic coding 的边界：
GPT-5.3-Codex 甚至参与了自身的训练调试和部署。
DeepSeek-V4 预计将于 2026 年 4 月发布，
以 mHC、Engram 等架构创新为基础，
结合原生多模态和 LTM 突破，
向闭源模型发起最强挑战。

2026 年的另一个标志性事件是\textbf{端侧模型的实用化}。
Step 3.5 Flash（196B 总参数，11B 活跃参数）
可以在消费级 RTX 4090 上本地运行，
性能接近 GPT-4 水平。
Qwen3.5-0.8B 仅 800M 参数，
可以在手机上运行推理，
用于离线翻译、日程管理和简单的代码辅助。
这标志着 LLM 从「云端独占」走向「端云协同」：
简单任务在本地处理（零延迟、无隐私泄露），
复杂任务上传到云端处理。

这场效率革命的意义远超技术本身。
当训练一个 frontier 模型的成本从数亿美元降至数百万美元，
AI 研究的参与门槛大幅降低。
大学实验室、初创企业甚至独立研究者
都有可能做出有意义的贡献。
这对整个领域的健康发展至关重要：
创新不应该被垄断在少数拥有万卡集群的机构手中。

这场效率革命还产生了一个深层的地缘政治影响。
美国对中国的芯片出口管制（2022 年 10 月起）
限制了中国获取最先进的 NVIDIA GPU。
但效率革命表明：\textbf{芯片算力的限制可以部分被算法效率弥补}。
DeepSeek-V3 使用的 H800 是被阉割过的 H100
（inter-chip 通信带宽减半），
但通过精心的工程优化（如 FP8 训练、自研通信库），
它实现了与使用全规格 H100 的西方团队相当的训练效率。
GLM-5 更进一步，完全使用华为 Ascend 910B 训练：
虽然单卡性能不如 H100，但通过 10 万卡的规模和软件优化，
仍然训练出了 frontier 级别的模型。
这证明了一个重要的论点：
在 AI 竞赛中，\textbf{算法创新和工程能力的价值至少等于硬件优势}。

\subsection{Scaling Hypothesis 之争}

贯穿这段历史的是一个核心假说：
\textbf{Scaling Hypothesis}：
持续增加模型规模、数据量和计算量，
就能持续获得更强的能力。
Kaplan 等人~\cite{kaplan2020scaling} 在 2020 年的研究
给出了精确的幂律关系（Scaling Laws），
为这一假说提供了实证基础。

但到了 2025--2026 年，这一假说面临深刻的重构。

\textbf{预训练 scaling 遭遇瓶颈}。
Ilya Sutskever 在 NeurIPS 2024 的演讲中明确指出：
「我们已知的预训练方式将会终结，因为我们只有一个互联网。」
高质量文本数据正在接近耗尽：
据估计，公开可用的高质量文本数据总量约为 10--15T tokens，
而 frontier 模型的训练数据量已逼近这一上限。
合成数据（synthetic data）虽然是一条出路，
但如何避免模型在自身生成的数据上训练导致的退化
（所谓「模型坍缩」）仍是开放问题。
一种有效的缓解策略是使用\textbf{任务导向的合成数据}：
不是让模型生成自由文本，
而是让模型解决特定的问题（如数学题、代码题），
然后用正确答案作为训练数据。
这种方式下，合成数据的质量可以通过验证器（如数学证明检查器、代码单元测试）来保证：
即使生成过程有噪声，只要最终答案通过了验证，数据就是高质量的。
DeepSeek-R1~\cite{deepseekr1} 的 RL 训练本质上就是这种思路的极致体现：
模型生成推理过程，环境验证答案正确性，正确的推理过程成为高质量训练信号。

\textbf{Test-time compute 开辟新维度}。
OpenAI 的 o 系列模型（o1/o3）和 DeepSeek-R1 证明了一条新路径：
不增大模型本身，而是让模型在推理时花更多计算来「思考」。
这本质上是将 scaling 从训练阶段转移到推理阶段。
2025 年 NeurIPS 的教程「Scale Test-Time Compute on Modern Hardware」
系统总结了这一范式：test-time compute 的 scaling law
与训练阶段有本质不同：
推理负载具有低并行度、不规则工作量和动态执行路径的特点，
在现代硬件上的高效部署是一个独立的技术挑战。

\textbf{从「规模时代」到「研究时代」}。
一个越来越强的共识是：AI 正在从「规模时代」（Age of Scaling）
过渡到「研究时代」（Age of Research）：
单纯增加 GPU 集群和数据量的策略回报递减，
突破需要来自架构创新（如 DeepSeek 的 MLA 和 Sparse Attention）、
训练方法创新（如纯 RL 训练推理能力）、
以及数据工程创新（如 agentic task synthesis pipeline）。
Meta 的 Llama~4 Behemoth（2T 参数）的延期发布
就是「更大不一定更好」的一个注脚：
内部评估发现它未必显著优于前代模型，
促使 Meta 转向更注重架构效率的策略。

但 scaling 并未「死亡」，它在\textbf{变形}。
DeepSeek-V3.2 通过 scalable RL 框架
在后训练阶段 scale compute 至 GPT-5 水平；
GPT-5 将推理和生成统一在一个模型中自动路由计算资源；
Qwen3 的混合推理模式让用户动态控制思考深度。
「规模」的含义已经从
「更大的模型」扩展到了一个多维空间：
更多的训练数据、更好的数据质量、
更高效的架构、更智能的后训练、
更多的推理时计算：
以及这些维度之间的最优分配。

\section{2026 年的格局}

截至 2026 年 3 月，LLM 领域已经形成了多极竞争的格局，
闭源与开源的边界日益模糊。

\subsection{闭源阵营}

\textbf{OpenAI} 在 2025 年完成了从 GPT 系列到统一模型的转型。
GPT-5（2025 年 8 月）将生成模型（GPT 系列）和推理模型（o 系列）
统一为单一系统，内置智能路由器自动分配计算资源。
后续的 GPT-5.1、5.2（2025 年秋）、5.3-Codex（2026 年 2 月）
和 GPT-5.4（2026 年 3 月）以每隔数周的节奏迭代，
GPT-5.2/5.4 系列在推理基准上持续刷新纪录，
GPT-5.3-Codex 甚至参与了自身的训练调试。
此前的 o3/o4-mini（2025 年 4 月）开创了推理模型的工具调用能力，
已在 GPT-5 发布后从 ChatGPT 主界面退役，但仍可通过 API 访问。

\textbf{Anthropic} 以 Claude 系列在 agentic coding 领域建立了强势地位。
Claude~4（2025 年 5 月）率先展示了长时间自主编程能力，
Claude Sonnet~4.5（2025 年 9 月）和 Opus~4.5（2025 年 11 月）持续领先，
2026 年 2 月的 \textbf{Claude Opus~4.6} 和 \textbf{Sonnet~4.6}
在 Terminal-Bench~2.0 和 Humanity's Last Exam 上取得最高分，
首次为 Opus 级模型提供百万 token 上下文窗口。
Claude Opus~4.6 在 SWE-bench Verified 上以 79.2\%（Thinking 模式）
刷新了 agentic coding 的纪录，
并引入了可配置的 thinking budget：
用户可精确控制推理计算的预算，
在通用模型中按需激活深度推理，而非训练独立的推理专用模型。

\textbf{Google} 以 Gemini~2.5~Pro（2025 年 3 月）重新加入竞争，
Gemini~3（2025 年 11 月）引入原生 thinking 模式和原生多模态，
\textbf{Gemini~3.1~Pro}（2026 年 2 月）
以原生思维链推理进一步提升了复杂推理和 agentic 能力，
以 1M token 上下文窗口覆盖代码库级别的分析任务。

\subsection{开源阵营的历史性追赶}

2025--2026 年，开源模型首次在质量上真正逼近甚至超越闭源模型。

\textbf{DeepSeek} 是这场追赶的先锋。
V3（2024 年 12 月）以 560 万美元达到 GPT-4 水平~\cite{deepseekv3}；
R1（2025 年 1 月）以纯 RL 训练达到 o1 水平~\cite{deepseekr1}；
V3.2（2025 年 12 月）通过 Sparse Attention 和 scalable RL
达到 GPT-5 水平，其高计算变体在 2025 年 IMO 和 IOI 中
取得金牌级表现。
备受期待的 \textbf{V4}（预计 2026 年 4 月）
将聚焦于原生多模态、LTM（Long-Term Memory）突破
以及国产芯片适配，
以约 1T 总参数和百万 token 上下文，
目标是成为开源模型对闭源的最强挑战。

\textbf{Meta} 的 Llama~4（2025 年 4 月）
发布了 Scout 和 Maverick 两款 MoE 多模态模型，
其中 Scout 支持高达 1000 万 token 的上下文窗口。
但旗舰级 Behemoth（2T 参数）因内部评估显示其
未必显著优于前代而反复延期，至今未发布，
成为「更大不一定更好」的典型案例。

\textbf{阿里巴巴} 的 Qwen3~\cite{qwen2025qwen3}（2025 年 4 月）
首创混合推理模式，旗舰 Qwen3-235B-A22B 在多项基准上与 DeepSeek-R1 持平。
2026 年 2 月发布的 \textbf{Qwen3.5} 以 397B 旗舰模型为核心，
采用 Gated DeltaNet 混合架构（3:1 线性/全注意力比例），
支持 201 种语言和 262K 上下文，
并推出 0.8B--9B 全系列小模型覆盖边缘场景：
其中 Qwen3.5-9B 在多项基准上超越前代 Qwen3-30B。

\textbf{2026 年春季的「模型发布海啸」}
是中国开源 AI 实力的集中展示。
从 1 月 27 日到 2 月 16 日，短短三周内，
6 家中国头部实验室密集发布了各自的旗舰模型：
Kimi K2.5（1T MoE，Agent Swarm 多智能体协作）、
Step 3.5 Flash（196B/11B，可在 RTX 4090 上本地运行）、
GLM-5（744B 参数，完全在华为 Ascend 910B 上训练，\textbf{零 NVIDIA 依赖}）、
MiniMax M2.5（SWE-bench Verified 80.2\%，超越同期所有闭源模型）、
Doubao Seed 2.0（AIME 2025 98.3\%）、
以及 Qwen 3.5（397B/17B）。
路透社将这一现象定性为「抢跑式竞争」（preemptive competition）：
各家在 DeepSeek V4 发布前争相卡位。
据 OpenRouter 统计，中国开源模型在全球开源 LLM 使用量中的占比
已从 2024 年的 1.2\% 攀升至 2026 年初的近 30\%，
标志着全球开源 AI 格局的根本性重构。

\subsection{关键趋势}

\textbf{推理模型的普及与融合}。
从 OpenAI o1/o3 到 DeepSeek-R1，
再到 GPT-5 的内置推理路由和 Claude 的 extended thinking，
「慢速思考」已从独立的推理模型范式
演变为通用模型的标准能力。
Qwen3 的 \texttt{/think} 和 \texttt{/no\_think} 切换，
以及 Gemini~3 的 Deep Think 模式，
都体现了同一趋势：用户可按需控制推理深度，
在响应速度与推理质量之间灵活权衡。

\textbf{多模态的原生化}。
当前的 frontier 模型普遍支持文本、图片、音频甚至视频的理解和生成。
Llama~4 和 Qwen3.5 从架构层面原生支持多模态，
不再是在文本模型上「外挂」视觉编码器。
纯文本 LLM 正在被多模态大模型取代：
但 Transformer 架构和预训练范式的核心不变。

\textbf{Agentic AI 的成熟}。
2026 年的模型不再仅仅是问答工具，
而是能够自主执行复杂工作流的 agent。
从 Claude Code 的多小时自主编程到 GPT-5.3-Codex 参与自身开发，
从 Gemini~3.1~Pro 的高级工具使用到 Qwen3.5 的自主 agent 能力，
AI 正在从「助手」演进为「同事」。
SWE-bench Verified 的分数从 2024 年初的约 5\%
跃升至 2026 年初的约 79\%（Claude Opus~4.6 Thinking），
Agent 能力在两年内实现了质的飞跃。
MCP（Model Context Protocol）已成为 Agent 工具交互的事实标准，
累计下载量突破 9700 万次，
由 Linux Foundation 接管治理，
OpenAI、Google、Microsoft 等主要玩家全部采纳。

\textbf{LTM（Long-Term Memory）崛起为 2026 年的「皇冠明珠」}。
传统的上下文窗口扩展正在被更本质的记忆架构取代。
DeepSeek 的 Engram 论文提出将「记忆」与「计算」分离，
通过 O(1) 查询的现代化 N-gram 嵌入实现可扩展的知识检索。
Google 的 Titans 架构引入了自更新的神经记忆模块，
能够在运行时持续学习，处理跨周/跨月的对话记忆。
LTM 有望成为下一代模型的核心差异化能力。

\textbf{国产芯片独立性的里程碑}。
GLM-5 在 10 万块华为 Ascend 910B 上完成了 744B 参数模型的完整训练，
实现了首个 frontier 级别模型的零 NVIDIA 依赖。
这标志着中国 AI 基础设施从「受限可用」迈向「独立可控」。

\textbf{「智能太便宜以至于无需计量」}。
MiniMax M2.5 的 API 定价仅为 GPT-5 的 1/10--1/20，
Qwen 3.5 比竞品便宜约 60\%。
推理成本的急剧下降正在改变整个产业的经济模型：
AI 能力的边际成本趋近于零的未来已不再遥远。

\textbf{评估基准的军备竞赛}。
传统基准（MMLU、GSM8K）在 frontier 模型上已趋于饱和。
2025--2026 年涌现了一批更具挑战性的评估：
SWE-bench Verified 测试模型解决真实软件工程问题的能力；
FrontierMath 包含研究级数学难题，即使 GPT-5 的通过率也低于 30\%；
Humanity's Last Exam 由各领域顶尖专家出题，
旨在测试 AI 与人类专家之间的差距；
Terminal-Bench 2.0 评估模型在终端环境中的自主操作能力。
这些新基准共同推动了模型能力的真实进步：
防止了「应试教育」式的基准刷分。

\textbf{开源与闭源的模糊边界}。
2025 年之前，「开源」和「闭源」是泾渭分明的阵营。
但到了 2026 年，这条边界变得模糊：
DeepSeek 完全开源模型权重和训练代码（MIT 许可证），
但不公开训练数据；
Meta 开放了 Llama 4 的权重，
但附带了使用限制（用户数超过 7 亿需要特殊许可）；
OpenAI 开放了 GPT-5 的 API 但不公开任何权重；
Anthropic 发表了详细的技术论文但不开放模型。
「开放程度」已经成为一个连续的光谱，
而非二元的分类。

\textbf{训练成本的戏剧性下降}。
LLM 训练成本的演变本身就是一条引人入胜的曲线：
\begin{center}\small
\renewcommand{\arraystretch}{1.3}
\begin{tabular}{lrr}
\toprule
\rowcolor{figLightGray}
\textbf{模型} & \textbf{年份} & \textbf{估算训练成本} \\
\midrule
GPT-3 (175B)      & 2020 & \$4.6M \\
Chinchilla (70B)  & 2022 & \$3--5M \\
PaLM (540B)       & 2022 & \$10--15M \\
GPT-4 (推测 $>$1T) & 2023 & \$50--100M \\
Llama 3 (405B)    & 2024 & \$30--50M \\
DeepSeek-V3 (671B)& 2024 & \$5.6M \\
GPT-5             & 2025 & \$100M+ \\
DeepSeek-V3.2     & 2025 & $\sim$\$10M \\
\bottomrule
\end{tabular}
\end{center}

这张表揭示了一个重要趋势：
同等能力水平上的训练成本在快速下降。
GPT-4 级别的性能从约 1 亿美元（GPT-4）
降至约 560 万美元（DeepSeek-V3）：
两个数量级的下降，仅用了不到两年。
这种下降不是来自硬件进步
（H800 和 A100 的性能差距远没有 20 倍），
而是来自\textbf{软件和方法论的创新}：
MoE 架构、FP8 训练、更高效的注意力机制、
更好的数据工程和训练策略。

这些成就背后有一个共同的故事：
更聪明的训练方法、更高效的架构设计、更精细的数据工程。
本文将完整地讲述这个故事：
从数据采集到模型部署，
从 Transformer 的数学原理到分布式训练的工程实践，
从预训练的计算规模定律到后训练的强化学习技巧。

我们的旅程从最基础的构件开始：Transformer 架构。

\section{本书的结构与阅读指南}

本书涵盖了 LLM 训练的完整技术栈，
从底层架构到高层应用，从理论基础到工程实践。
以下是各章节的逻辑关系和阅读建议。

\textbf{基础层}（第 2--3 章）。
第~\ref{ch:transformer}章详解 Transformer 架构：
注意力机制、位置编码、归一化方案、激活函数的设计选择及其背后的动机。
第~\ref{ch:tokenization}章讨论分词：
一个看似简单但对模型能力有深远影响的关键组件。
这两章是理解后续内容的必要基础。

\textbf{数据与训练}（第 4--5 章）。
第~\ref{ch:data}章覆盖数据工程：
数据的采集、清洗、去重、配比，
以及合成数据和数据质量评估。
第~\ref{ch:pretrain}章讨论预训练的工程实践：
分布式并行化策略、混合精度训练、训练稳定性。
这两章是理解「如何高效地训练一个大模型」的核心。

\textbf{架构创新}（第 6--7 章）。
第~\ref{ch:moe}章深入 MoE（混合专家）架构：
DeepSeek-V3 和 Llama 4 的核心技术，
理解「如何用更少的计算获得更强的能力」。
第~\ref{ch:long-context}章讨论长上下文处理：
从位置编码的扩展到注意力机制的稀疏化，
理解模型如何处理百万 token 级别的输入。

\textbf{后训练与推理}（第 8--10 章）。
第~\ref{ch:post-training}章覆盖后训练：
SFT、RLHF、DPO 等对齐技术，
以及 DeepSeek-R1 的纯 RL 范式。
第~\ref{ch:inference}章讨论推理优化：
量化、蒸馏、投机解码等部署技术。
第~\ref{ch:ttc}章聚焦 test-time compute：
推理模型的 scaling 策略。

\textbf{评估与前沿}（第 11--13 章）。
第~\ref{ch:eval}章讨论评估方法论，
第~\ref{ch:org}章分析各大研究团队的技术路线，
第~\ref{ch:outlook}章展望未来方向。

每一章都可以独立阅读，但我们建议按顺序通读基础层（第 2--3 章），
然后根据兴趣选读后续章节。

\subsection{本章的核心论点回顾}

在进入技术细节之前，让我们回顾本章建立的核心论点：

\textbf{论点一：Next-token prediction 是通向通用智能的有效路径}。
这不是因为「预测下一个词」这个任务本身有什么特别：
它的特别之处在于\textbf{通用性}：
要在所有可能的上下文中做出正确预测，
模型必须内化人类知识的方方面面。
压缩-预测-理解的等价性为这个直觉提供了信息论的基础。

\textbf{论点二：规模是能力的关键驱动力，但方式在变化}。
从 Kaplan 到 Chinchilla，从预训练 scaling 到 test-time compute，
「规模」的含义已经从单纯的「更多参数」
扩展为一个多维的资源分配问题。
2025--2026 年的效率革命证明了：
更聪明的架构和训练方法可以大幅降低达到同等能力所需的计算量。

\textbf{论点三：涌现是规模的产物，但机制尚不完全清楚}。
大模型能做到小模型做不到的事——这是经验事实。
涌现是否是「真实的相变」还是「度量假象」仍有争议，
但从应用角度来说，这个区分不太重要：
无论底层机制如何，增大规模确实能解锁新能力。

\textbf{论点四：LLM 训练是一项系统工程}。
从数据清洗到分词器设计，从分布式训练到后训练对齐，
每个环节都有大量的技术细节和设计选择。
没有任何一个单独的技巧可以「一劳永逸」：
frontier 模型的成功是数十项工程决策的累积效果。

这些论点将在后续章节中被反复印证和深化。
最后，一条实用的阅读建议。
本书的每一章都包含大量的技术细节：
数学公式、架构图、代码片段、训练技巧。
对于初次阅读的读者，
不必追求每一个细节都理解：
先建立全局概念框架，
然后在实际需要时回来查阅具体细节。
对于有实践经验的读者，
我们希望本书能为你已有的知识提供系统化的组织，
以及对最新进展的全面覆盖。

LLM 的发展速度之快令人叹为观止。
在写作本书的过程中，
我们多次需要更新「最新」的数据和模型：
因为每隔几周就有新的 frontier 模型发布。
本书反映的是截至 2026 年 3 月的技术状态。
读者阅读时，某些具体的数字和排名可能已经过时，
但底层的原理和方法论：
Transformer 的注意力机制、Scaling Laws 的幂律关系、
RLHF 的偏好优化、BPE 的子词切分：
这些在可预见的未来不会改变。

现在，让我们正式开始技术旅程。


%% ============================================================
%% NEW REFERENCES (to be added to refs.bib):
%% shannon1950 = Shannon, Claude E. "Prediction and Entropy of Printed English." Bell System Technical Journal, 30(1):50--64, 1951 (presented 1950).
%% deletang2024language = Delétang, Grégoire et al. "Language Modeling Is Compression." arXiv:2309.10668, 2023 (published ICLR 2024).
%% wei2022cot = Wei, Jason et al. "Chain-of-Thought Prompting Elicits Reasoning in Large Language Models." NeurIPS 2022.
%% schaeffer2023mirage = Schaeffer, Rylan et al. "Are Emergent Abilities of Large Language Models a Mirage?" NeurIPS 2023.
%% mikolov2013word2vec = Mikolov, Tomas et al. "Efficient Estimation of Word Representations in Vector Space." arXiv:1301.3781, 2013.
%% peters2018elmo = Peters, Matthew E. et al. "Deep Contextualized Word Representations." NAACL 2018.
%% radford2019gpt2 = Radford, Alec et al. "Language Models are Unsupervised Multitask Learners." OpenAI blog, 2019.
%% yang2019xlnet = Yang, Zhilin et al. "XLNet: Generalized Autoregressive Pretraining for Language Understanding." NeurIPS 2019.
%% deng2009imagenet = Deng, Jia et al. "ImageNet: A Large-Scale Hierarchical Image Database." CVPR 2009.
%% sardana2024chinchilla = Sardana, Nikhil and Frankle, Jonathan. "Beyond Chinchilla-Optimal: Accounting for Inference in Language Model Scaling Laws." ICML 2024.
%% penedo2024fineweb = Penedo, Guilherme et al. "The FineWeb Datasets: Decanting the Web for the Finest Text Data at Scale." NeurIPS 2024 Datasets Track.
%%
%% hahn2024theory = Hahn, Michael and Goyal, Navin. "A Theory of Emergent In-Context Learning as Implicit Structure Induction." arXiv:2303.07971, 2023.
%% li2023othello = Li, Kenneth et al. "Emergent World Representations: Exploring a Sequence Model Trained on a Synthetic Task." ICLR 2023.
%% olsson2022induction = Olsson, Catherine et al. "In-context Learning and Induction Heads." Transformer Circuits Thread, 2022.
%% arora2023theory = Arora, Sanjeev and Goyal, Anirudh. "A Theory for Emergence of Complex Skills in Language Models." arXiv:2307.15936, 2023.
%% wei2023larger = Wei, Jason et al. "Larger language models do in-context learning differently." arXiv:2303.03846, 2023.
%% pearl2009causality = Pearl, Judea. "Causality: Models, Reasoning, and Inference." Cambridge University Press, 2009.
%% merrill2023expressive = Merrill, William and Sabharwal, Ashish. "The Expressive Power of Transformers with Chain of Thought." ICLR 2024.
%% rafailov2023dpo = Rafailov, Rafael et al. "Direct Preference Optimization: Your Language Model is Secretly a Reward Model." NeurIPS 2023.
%% glazer2024frontiermath = Glazer, Elliot et al. "FrontierMath: A Benchmark for Evaluating Advanced Mathematical Reasoning in AI." arXiv:2411.04872, 2024.
%% rajbhandari2020zero = Rajbhandari, Samyam et al. "ZeRO: Memory Optimizations Toward Training Trillion Parameter Models." SC, 2020.
%% li2023othello = Li, Kenneth et al. "Emergent World Representations: Exploring a Sequence Model Trained on a Synthetic Task." ICLR, 2023.
%%
%% --- Added 2026-03-08 ---
%% deepseekv32 = DeepSeek-AI. "DeepSeek-V3.2: Pushing the Frontier of Open Large Language Models." arXiv:2512.02556, 2025.
%% qwen2025qwen35 = Qwen Team. "Qwen3.5 Technical Report." Alibaba Cloud, 2026.
%% openai2025gpt5 = OpenAI. "Introducing GPT-5." openai.com, August 2025.
%% openai2025o3 = OpenAI. "Introducing OpenAI o3 and o4-mini." openai.com, April 2025.
%% anthropic2025claude4 = Anthropic. "Introducing Claude 4." anthropic.com, May 2025.
%% anthropic2026opus46 = Anthropic. "Introducing Claude Opus 4.6." anthropic.com, February 2026.
%% google2025gemini3 = Google DeepMind. "A New Era of Intelligence with Gemini 3." blog.google, November 2025.
%% google2026gemini31 = Google DeepMind. "Gemini 3.1 Pro: A Smarter Model for Your Most Complex Tasks." blog.google, February 2026.
%% meta2025llama4 = Meta AI. "The Llama 4 Herd." ai.meta.com, April 2025.
%% ============================================================
